%% Lecture 18 -- Effective magnetic charges, linear magnetic media, uniqueness, and forces
\section{Lecture 18 -- Effective Magnetic Charges, Linear Media, and Forces on Magnetic Matter}\label{sec:lec18}

\begin{center}
\begin{tabular}{ll}
  \textbf{Date:}\quad 07/06/2026 & \\
\end{tabular}
\end{center}
\hrule\vspace{1.5em}

\subsection*{Overview}

Lecture 17 reduced a magnetized body to a magnetic scalar potential $\psi_M$ obeying the
electrostatic Poisson equation $\lapl\psi_M=-4\pi\rho_M$ with $\rho_M=-\divg\bvec{M}$.
This lecture turns that formal analogy into a working method and then studies the
materials themselves:
\begin{itemize}
  \item complete the proof of the curl form of Gauss's theorem used last time;
  \item rewrite $\psi_M$ as the potential of \emph{effective volume and surface magnetic
        charges} $\rho_M=-\divg\bvec{M}$, $\sigma_M=\bvec{M}\cdot\uvec{n}$, and use them
        as a recipe to solve a uniformly magnetized magnet;
  \item state the matching conditions on $\bB$ and $\Hvec$ across an interface;
  \item introduce \emph{linear media} $\bvec{M}=\chi\Hvec$ with permeability
        $\mu=1+4\pi\chi$ (paramagnets $\chi>0$, diamagnets $\chi<0$);
  \item prove \emph{uniqueness} of the resulting boundary value problem and deduce that
        linear media carry no remanent magnetization;
  \item compute the \emph{force} on linear matter, $\propto\nabla  B^{2}$, explaining why
        paramagnets are pulled into strong fields and diamagnets expelled;
  \item begin the microscopic (Langevin) account of diamagnetism.
\end{itemize}

%------------------------------------------------------------
\subsection{The Curl Form of Gauss's Theorem}\label{subsec:lec18-curl-gauss}

\subsubsection*{Mathematical Setup}
In Lecture 17 we converted the curl term of the magnetization vector potential into a
surface integral using the identity
\begin{equation}\label{eq:lec18-curl-gauss}
  \int_V (\nabla\times\bvec{A})\,d^3r=\oint_S \uvec{n}\times\bvec{A}\,d\sigma,
\end{equation}
with $\uvec{n}$ the outward normal of $S=\partial V$. It looks unlike the ordinary
divergence theorem; here we prove it by reducing it to the standard one. The trick is to
dot \eqref{eq:lec18-curl-gauss} into an \emph{arbitrary constant} vector $\bvec{c}$ and
to use a vector identity.

\lem[lem:lec18-curl-gauss]{Curl Form of the Divergence Theorem}{
  For any smooth vector field $\bvec{A}$ on a volume $V$ with boundary $S$ and outward
  normal $\uvec{n}$,
  \begin{equation}\label{eq:lec18-curl-gauss-boxed}
    \boxed{\int_V (\nabla\times\bvec{A})\,d^3r=\oint_S \uvec{n}\times\bvec{A}\,d\sigma.}
  \end{equation}
}

\subsubsection*{Proof}
Let $\bvec{c}$ be a constant vector and dot it into the left-hand side. Since $\bvec{c}$
is constant it passes inside the integral:
\begin{equation}\label{eq:lec18-cdot-curl}
  \bvec{c}\cdot\!\int_V (\nabla\times\bvec{A})\,d^3r
  =\int_V \bvec{c}\cdot(\nabla\times\bvec{A})\,d^3r.
\end{equation}
Use the product identity (an $\varepsilon_{ijk}$ one-liner)
\begin{equation}\label{eq:lec18-vec-identity}
  \divg(\bvec{A}\times\bvec{c})
  =\bvec{c}\cdot(\nabla\times\bvec{A})-\bvec{A}\cdot(\nabla\times\bvec{c}).
\end{equation}
Because $\bvec{c}$ is constant, $\nabla\times\bvec{c}=0$ and the last term drops, so
$\bvec{c}\cdot(\nabla\times\bvec{A})=\divg(\bvec{A}\times\bvec{c})$. Therefore
\begin{equation}\label{eq:lec18-apply-divthm}
  \int_V \bvec{c}\cdot(\nabla\times\bvec{A})\,d^3r
  =\int_V \divg(\bvec{A}\times\bvec{c})\,d^3r
  =\oint_S (\bvec{A}\times\bvec{c})\cdot\uvec{n}\,d\sigma,
\end{equation}
the last step being the ordinary divergence theorem applied to the vector
$\bvec{A}\times\bvec{c}$. The integrand is a scalar triple product; cycling it,
\begin{equation}\label{eq:lec18-cyclic}
  (\bvec{A}\times\bvec{c})\cdot\uvec{n}
  =\bvec{c}\cdot(\uvec{n}\times\bvec{A}),
\end{equation}
so that
\begin{equation}\label{eq:lec18-cancel-c}
  \bvec{c}\cdot\!\int_V (\nabla\times\bvec{A})\,d^3r
  =\bvec{c}\cdot\!\oint_S (\uvec{n}\times\bvec{A})\,d\sigma.
\end{equation}
Since $\bvec{c}$ is an \emph{arbitrary} constant vector, the two vectors it multiplies
must be equal, which is \eqref{eq:lec18-curl-gauss-boxed}. $\blacksquare$

\nt{The pattern ``dot into a constant $\bvec{c}$, reduce to the standard divergence
theorem, then strip $\bvec{c}$'' generates a whole family of integral theorems
($\int\nabla \phi\,d^3r=\oint\phi\,\uvec{n}\,d\sigma$, etc.) from the single divergence
theorem.}

%------------------------------------------------------------
\subsection{Effective Magnetic Charges}\label{subsec:lec18-magnetic-charges}

\subsubsection*{Mathematical Setup}
Start from the magnetic scalar potential of Lecture 17,
\begin{equation}\label{eq:lec18-psi-start}
  \psi_M(\br)=\int d^3r'\,
  \frac{\bvec{M}(\br')\cdot(\br-\br')}{|\br-\br'|^3}
  =\int d^3r'\,\bvec{M}(\br')\cdot\nabla'\frac{1}{|\br-\br'|},
\end{equation}
where in the second form we used
$\nabla'|\br-\br'|^{-1}=(\br-\br')/|\br-\br'|^3$ (the primed gradient acts on $\br'$).
Now run the same product-rule manipulation as for the bound currents, but with a
\emph{divergence} instead of a curl. Write
\begin{equation}\label{eq:lec18-product}
  \divg'\!\left(\frac{\bvec{M}}{|\br-\br'|}\right)
  =\bvec{M}\cdot\nabla'\frac{1}{|\br-\br'|}
  +\frac{\divg'\bvec{M}}{|\br-\br'|},
\end{equation}
so the integrand of \eqref{eq:lec18-psi-start} becomes
\begin{equation}\label{eq:lec18-integrand}
  \bvec{M}\cdot\nabla'\frac{1}{|\br-\br'|}
  =\divg'\!\left(\frac{\bvec{M}}{|\br-\br'|}\right)
  -\frac{\divg'\bvec{M}}{|\br-\br'|}.
\end{equation}
Integrating over the body and turning the first (total-divergence) term into a surface
integral by the ordinary divergence theorem,
\begin{equation}\label{eq:lec18-psi-split}
  \boxed{\psi_M(\br)
  =\oint_S \frac{\bvec{M}(\br')\cdot\uvec{n}}{|\br-\br'|}\,d\sigma'
  +\int_V \frac{-\divg'\bvec{M}(\br')}{|\br-\br'|}\,d^3r'.}
\end{equation}
\textit{Significance:} $\psi_M$ is literally the electrostatic potential of a surface
charge $\sigma_M=\bvec{M}\cdot\uvec{n}$ plus a volume charge $\rho_M=-\divg\bvec{M}$; the
magnetic problem is now an electrostatics problem.

\dfn[dfn:lec18-eff-charges]{Effective Magnetic Charges}{
  In a region with no free currents the magnetized body is equivalent, for the purpose
  of computing $\Hvec=-\nabla \psi_M$, to the fictitious charge densities
  \begin{equation}\label{eq:lec18-eff-charges}
    \boxed{\rho_M=-\divg\bvec{M}\quad(\text{volume}),
    \qquad
    \sigma_M=\bvec{M}\cdot\uvec{n}\quad(\text{surface}).}
  \end{equation}
}
\textit{Significance:} These are bookkeeping sources only -- there are no magnetic
monopoles -- but every electrostatic technique (images, separation of variables,
multipole expansion) transfers directly to magnetostatics through them. They are the
magnetic mirror of the bound charges $\rho_b=-\divg\bvec{P}$, $\sigma_b=\bvec{P}\cdot\uvec{n}$
of a dielectric.

\nt{The same physics also has a \emph{current} description from Lecture 17,
$\Jvec_M=c\,\nabla\times\bvec{M}$ and $\bvec{K}_M=c\,\bvec{M}\times\uvec{n}$. The two
pictures (effective charges vs.\ bound currents) are completely equivalent; which is
easier is a matter of symmetry and convenience, exactly like choosing a coordinate
system.}

\subsubsection*{Worked Recipe: the Uniformly Magnetized Magnet}
Consider a permanent magnet -- take a cylinder of height $h$ and radius $a$ -- with
\emph{uniform} magnetization $\bvec{M}=M\uvec{z}$ along its axis, sitting in vacuum with
no free currents. We want $\bB$ everywhere.

\ex[ex:lec18-bar-magnet]{Uniformly Magnetized Cylinder}{
  Because $\bvec{M}$ is constant, $\rho_M=-\divg\bvec{M}=0$: there is \emph{no} volume
  charge. The only sources are surface charges $\sigma_M=\bvec{M}\cdot\uvec{n}=M\,(\uvec{z}\cdot\uvec{n})$:
  \begin{itemize}
    \item \textbf{top cap} ($\uvec{n}=+\uvec{z}$): $\sigma_M=+M$;
    \item \textbf{bottom cap} ($\uvec{n}=-\uvec{z}$): $\sigma_M=-M$;
    \item \textbf{side wall} ($\uvec{n}\perp\uvec{z}$): $\sigma_M=0$.
  \end{itemize}
  Solve the \emph{electrostatics} of these two charged disks for the potential
  $\psi_M(\br)$ (an ordinary $\sigma/|\br-\br'|$ integral), then
  \begin{equation}\label{eq:lec18-recipe}
    \Hvec=-\nabla \psi_M,
    \qquad
    \bB=\Hvec+4\pi\bvec{M}.
  \end{equation}
  Outside the magnet $\bvec{M}=0$, so $\bB=\Hvec$; inside they differ by the uniform
  $4\pi\bvec{M}$. Because $\bvec{M}$ is constant the addition is trivial, and the field
  is obtained in all of space.
}
\textit{Significance:} A magnet with uniform $\bvec{M}$ produces exactly the $\Hvec$-field
of two oppositely charged disks -- the magnetic ``north'' and ``south'' poles are the
$\pm\sigma_M$ caps.

\begin{figure}[h]
\centering
\begin{tikzpicture}[>=Stealth, thick, scale=1.0]
  % the magnet body
  \draw[rounded corners=1pt, fill=blue!6] (-0.9,-1.6) rectangle (0.9,1.6);
  % magnetization arrow
  \draw[->, red!70!black, line width=1.2pt] (0,-1.0) -- (0,1.0)
    node[midway, right] {$\bvec{M}=M\uvec{z}$};
  % top positive charges
  \node at (-0.55,1.35) {$+$}; \node at (-0.18,1.35) {$+$};
  \node at (0.18,1.35) {$+$};  \node at (0.55,1.35) {$+$};
  \node[above] at (0,1.6) {$\sigma_M=+M$};
  % bottom negative charges
  \node at (-0.55,-1.35) {$-$}; \node at (-0.18,-1.35) {$-$};
  \node at (0.18,-1.35) {$-$};  \node at (0.55,-1.35) {$-$};
  \node[below] at (0,-1.6) {$\sigma_M=-M$};
  % side label
  \node[right, gray] at (0.9,0) {$\sigma_M=0$};
  % axis
  \draw[->, gray] (1.9,-1.6) -- (1.9,1.7) node[above] {$z$};
\end{tikzpicture}
\caption{A uniformly magnetized cylinder. Since $\divg\bvec{M}=0$ there is no volume
magnetic charge; the equivalent electrostatic problem has surface charge
$\sigma_M=\bvec{M}\cdot\uvec{n}=\pm M$ only on the caps and zero on the side wall.}
\label{fig:lec18-bar-magnet}
\end{figure}

%------------------------------------------------------------
\subsection{Boundary Conditions for $\bB$ and $\Hvec$}\label{subsec:lec18-bc}

\subsubsection*{Conceptual Background}
When a problem has two regions (two materials, or matter and vacuum) one solves the
field equations in each region and stitches the solutions together at the interface. The
matching conditions follow from the two field laws exactly as in electrostatics.

\paragraph{Normal $\bB$ is continuous.}
From $\divg\bB=0$ (exact -- no magnetic monopoles) apply the divergence theorem to a
thin Gaussian pillbox straddling the interface; as its height $\to0$ the flux through the
rim vanishes and only the two faces survive, forcing
\begin{equation}\label{eq:lec18-Bperp}
  \boxed{B^{\perp}_{1}=B^{\perp}_{2}.}
\end{equation}
\textit{Significance:} The component of $\bB$ normal to any interface never jumps,
regardless of currents.

\paragraph{Tangential $\Hvec$ is continuous (absent free surface currents).}
From $\nabla\times\Hvec=(4\pi/c)\Jvec_{\rm f}$, run a thin Stokesian loop across the
interface. The circulation equals the enclosed free current. A \emph{volume} free
current $\Jvec_{\rm f}$ gives a flux that vanishes as the loop is squeezed onto the
surface; only a \emph{free surface current} $\bvec{K}_{\rm f}$ survives. Hence, if
$\bvec{K}_{\rm f}=0$,
\begin{equation}\label{eq:lec18-Hpar}
  \boxed{H^{\parallel}_{1}=H^{\parallel}_{2}\qquad(\bvec{K}_{\rm f}=0).}
\end{equation}
\textit{Significance:} The tangential $\Hvec$ jumps \emph{only} where a genuine free
sheet current flows on the surface -- the magnetic analogue of how the tangential $\bE$
(here the normal field) jumps only at a surface charge. Bound/volume currents never cause
a jump because their enclosed flux shrinks with the loop.

\nt{There is no boundary condition on $\divg\Hvec$ and no general jump condition built
from it: we have a formula for $\nabla\times\Hvec$ but the ``source'' $-4\pi\divg\bvec{M}$
of $\Hvec$ is internal to the matter. Use $B^{\perp}$ continuous and $H^{\parallel}$
continuous as the working pair.}

%------------------------------------------------------------
\subsection{Linear Magnetic Media}\label{subsec:lec18-linear}

\subsubsection*{Physical Motivation}
So far $\bvec{M}$ was \emph{given}. The more common physical situation is the opposite:
an initially unmagnetized material is placed in an external field, it magnetizes, and the
induced $\bvec{M}$ in turn modifies the field around it. In general $\bvec{M}=\bvec{M}(\bB)$
is a self-consistent unknown, and the system
\begin{equation}\label{eq:lec18-system}
  \divg\bB=0,\qquad
  \nabla\times\Hvec=\frac{4\pi}{c}\Jvec_{\rm f},\qquad
  \bB=\Hvec+4\pi\bvec{M}
\end{equation}
is not closed until the constitutive relation $\bvec{M}(\Hvec)$ is supplied.

\dfn[dfn:lec18-linear]{Linear Media: Susceptibility and Permeability}{
  A \emph{linear} material (paramagnet or diamagnet) obeys
  \begin{equation}\label{eq:lec18-MchiH}
    \boxed{\bvec{M}=\chi\,\Hvec,}
  \end{equation}
  with $\chi$ the \emph{magnetic susceptibility} (an intrinsic material constant, possibly
  position-dependent in an inhomogeneous body). Substituting into $\bB=\Hvec+4\pi\bvec{M}$,
  \begin{equation}\label{eq:lec18-BmuH}
    \boxed{\bB=(1+4\pi\chi)\Hvec\equiv\mu\,\Hvec,
    \qquad
    \mu=1+4\pi\chi,}
  \end{equation}
  where $\mu$ is the \emph{magnetic permeability}.
}
\textit{Significance:} The single constant $\mu$ closes the system
\eqref{eq:lec18-system}; the whole response of the material is encoded in it.

The sign of $\chi$ classifies the material:
\begin{center}
\renewcommand{\arraystretch}{1.3}
\begin{tabular}{lll}
  \toprule
  $\chi>0$ & paramagnet & $\mu>1$ \\
  $\chi<0$ & diamagnet & $\mu<1$ \\
  $\chi=0$ & vacuum & $\mu=1$ \\
  \bottomrule
\end{tabular}
\end{center}

\nt{Every material responds to a magnetic field -- there is no truly ``non-magnetic''
matter. Everyday ``non-magnetic'' objects are simply dia/paramagnets with extremely weak
$\chi$: for diamagnets $\chi\sim10^{-10}$ (order of magnitude), so $\mu$ sits
extraordinarily close to $1$. Gold, for instance, is diamagnetic but its response is
invisible to ordinary magnets. Linearity itself is special: ferromagnets are nonlinear
and hysteretic, and the linear law is at best a weak-field approximation there.}

%------------------------------------------------------------
\subsection{Uniqueness of the Boundary Value Problem}\label{subsec:lec18-uniqueness}

\subsubsection*{Mathematical Setup}
With $\mu$ given, the system \eqref{eq:lec18-system} is closed. We show its solution is
\emph{unique}: given the free currents and the material data, there is at most one field.
Suppose two solutions $(\bB_1,\Hvec_1)$ and $(\bB_2,\Hvec_2)$ exist and form the
difference
\begin{equation}\label{eq:lec18-diff}
  \bB'=\bB_1-\bB_2,\qquad \Hvec'=\Hvec_1-\Hvec_2.
\end{equation}
Subtracting the two copies of \eqref{eq:lec18-system}, and using
$\bB_i=\mu\Hvec_i$, the difference satisfies the \emph{source-free} system
\begin{equation}\label{eq:lec18-diff-system}
  \divg\bB'=0,\qquad \nabla\times\Hvec'=0,\qquad \bB'=\mu\,\Hvec'.
\end{equation}

\subsubsection*{The Energy Argument}
Since $\nabla\times\Hvec'=0$, there is a potential $\Hvec'=-\nabla \psi'$. Consider the
positive-looking quantity $\mu\,|\Hvec'|^2=\bB'\cdot\Hvec'$ and rewrite it as a
divergence. Using $\Hvec'=-\nabla \psi'$,
\begin{equation}\label{eq:lec18-BdotH}
  \bB'\cdot\Hvec'=-\bB'\cdot\nabla \psi'
  =-\divg(\psi'\bB')+\psi'\,\divg\bB'
  =-\divg(\psi'\bB'),
\end{equation}
because $\divg\bB'=0$. Integrate over all space and convert to a surface integral over a
large sphere $S_R$ of radius $R$:
\begin{equation}\label{eq:lec18-integral}
  \int_{\text{all space}}\mu\,|\Hvec'|^2\,d^3r
  =-\oint_{S_R}\psi'\,\bB'\cdot\uvec{n}\,d\sigma.
\end{equation}
For sources confined to a finite region the fields fall off at least as fast as those of
a dipole: $\psi'\sim 1/R^2$ and $B'\sim1/R^3$ (even a hypothetical monopole would give
$\psi'\sim1/R$, $B'\sim1/R^2$, i.e.\ integrand $\sim1/R^3$). Since the area grows only as
$R^2$, the surface term vanishes as $R\to\infty$:
\begin{equation}\label{eq:lec18-zero}
  \boxed{\int_{\text{all space}}\mu\,|\Hvec'|^2\,d^3r=0.}
\end{equation}
\textit{Significance:} The total ``energy'' of the difference field is forced to zero by
the boundary behaviour alone.

\subsubsection*{Conclusion}
For the materials we treat $\mu>0$ everywhere (paramagnets have $\chi>0$ so $\mu>1$;
diamagnets have tiny $|\chi|$ so $\mu\lesssim1$ stays positive). Then the integrand
$\mu|\Hvec'|^2\ge0$, and a non-negative integrand with zero integral must vanish
pointwise:
\begin{equation}\label{eq:lec18-Hprime-zero}
  \Hvec'=0\ \Rightarrow\ \bB'=\mu\Hvec'=0
  \ \Rightarrow\ (\bB_1,\Hvec_1)=(\bB_2,\Hvec_2).
\end{equation}
The solution is unique. $\blacksquare$

\begin{remark}
The argument needs $\mu>0$, which is why it covers ordinary matter:
$\mu=1+4\pi\chi$ with $|\chi|\ll1$ sits just above (paramagnets) or just below
(diamagnets) unity. Nothing in the magnetostatic equations \emph{forbids} a strongly
negative $\chi$; the bound $\chi>-1/4\pi$ (so $\mu>0$) is a \emph{thermodynamic}
stability requirement, not an electromagnetic one. The extreme case
$\chi=-\tfrac{1}{4\pi}$ gives $\mu=0$, i.e.\ $\bB=\mu\Hvec=0$ inside the material: this is
the \emph{ideal diamagnet} -- a superconductor expelling all flux (the Meissner effect of
Lecture 16). Genuine $\mu<0$ media are exotic and unstable, so for every material we treat
$\mu>0$ and uniqueness holds.
\end{remark}

\cor[cor:lec18-no-remanence]{No Remanent Magnetization in Linear Media}{
  With no free currents ($\Jvec_{\rm f}=0$) the full system \eqref{eq:lec18-system}
  coincides with the source-free difference system \eqref{eq:lec18-diff-system}, whose
  unique solution is $\bB=\Hvec=0$, hence $\bvec{M}=\chi\Hvec=0$. A linear material left
  to itself carries \emph{no} field.
}
\textit{Significance:} Magnetization of a linear medium exists only while an external
source is on. Switch the free currents on, they induce $\bvec{M}$; switch them off and
$\bvec{M}$ disappears. This is exactly what distinguishes linear media from ferromagnets,
where the $\bvec{M}(\Hvec)$ relation is different and magnetization \emph{remains} after
the external field is removed (permanent magnets).

%------------------------------------------------------------
\subsection{Force on Linear Magnetic Matter}\label{subsec:lec18-force}

\subsubsection*{Physical Motivation}
A magnetic dipole $\bmu$ in a field has potential energy $U=-\bmu\cdot\bB$ and feels a
force $\bvec{F}=-\nabla  U=\nabla (\bmu\cdot\bB)$ (for $\bmu$ held fixed). For a piece of
linear matter the moment is itself induced by the field, $\bmu\propto\bB$, and the energy
acquires a compensating factor of $\tfrac12$ -- exactly as the electrostatic energy
$\tfrac12\int\rho\phi$ carries a $\tfrac12$ to avoid double-counting the
self-interaction.

\subsubsection*{Force Density}
Work with the force per unit volume. Using $\bvec{M}=\chi\Hvec$ and $\Hvec=\bB/\mu$, the
$i$-th component of the force density is
\begin{equation}\label{eq:lec18-fdensity-step}
  f_i=M_j\,\partial_i B_j
  =\frac{\chi}{\mu}\,B_j\,\partial_i B_j
  =\frac{\chi}{2\mu}\,\partial_i\!\left(B_j B_j\right)
  =\frac{\chi}{2\mu}\,\partial_i B^{2}.
\end{equation}
Here $B^2=\bB\cdot\bB$ is the \emph{total} field at the material element (this is where
the $\tfrac12$ of the energy argument is automatically reproduced: writing the induced
moment through $\bB$ counts the same energy twice, which the $\tfrac12$ corrects).
In vector form, for constant $\chi,\mu$,
\begin{equation}\label{eq:lec18-force-boxed}
  \boxed{\bvec{f}=\nabla \!\left(\frac{\chi}{2\mu}\,B^{2}\right)
  =-\nabla  U,\qquad
  U=-\frac{\chi}{2\mu}\,B^{2}.}
\end{equation}
\textit{Significance:} The force on linear matter is a gradient of $B^2$; its sign is set
by $\chi$, so paramagnets ($\chi>0$) and diamagnets ($\chi<0$) move oppositely. A uniform
field exerts \emph{no} net force ($\nabla  B^2=0$).

\subsubsection*{Para- vs.\ Diamagnets}
Because $\nabla  B^2$ points toward stronger field:
\begin{itemize}
  \item \textbf{Paramagnet} ($\chi>0$): $\bvec{f}$ points up the gradient -- the sample is
        pulled \emph{into} regions of strong field, i.e.\ it sticks to the pole of a
        magnet.
  \item \textbf{Diamagnet} ($\chi<0$): $\bvec{f}$ points down the gradient -- the sample
        is pushed \emph{out} toward weak field, i.e.\ it is repelled by a magnet.
\end{itemize}

\begin{figure}[h]
\centering
\begin{tikzpicture}[>=Stealth, thick, scale=1.0]
  % field gradient region: dense lines on the left (strong), sparse on the right (weak)
  \foreach \x/\op in {0/100, 0.5/80, 1.1/55, 1.8/35, 2.7/20} {
    \draw[blue!\op!gray] (\x,-1.3) -- (\x,1.3);
  }
  \node[blue!70!black] at (0.2,1.6) {strong $B$};
  \node[gray] at (2.7,1.6) {weak $B$};
  % gradient direction
  \draw[->, thick] (3.4,0) -- (4.2,0) node[right] {$\nabla  B^2$};
  % paramagnet pulled toward strong field
  \filldraw[red!25] (1.9,0.5) circle (0.18);
  \draw[->, red!70!black, line width=1pt] (1.7,0.5) -- (1.0,0.5);
  \node[red!70!black, right] at (1.95,0.5) {para ($\chi>0$)};
  % diamagnet pushed toward weak field
  \filldraw[green!25] (1.0,-0.6) circle (0.18);
  \draw[->, green!50!black, line width=1pt] (1.2,-0.6) -- (1.9,-0.6);
  \node[green!50!black, right] at (1.95,-0.6) {dia ($\chi<0$)};
\end{tikzpicture}
\caption{In an inhomogeneous field a paramagnet is drawn toward strong field (up the
$\nabla  B^2$) while a diamagnet is expelled toward weak field. A uniform field produces no
net force.}
\label{fig:lec18-para-dia}
\end{figure}

\nt{The everyday analogue is electrostatic: a scrap of paper (a dielectric, effective
$\chi>0$, ``like a paramagnet'') is drawn toward the strong field near a statically
charged garment, because the field induces dipoles and the polarized scrap is pulled into
the strong-field region -- the same $\nabla (\text{field})^2$ mechanism.}

%------------------------------------------------------------
\subsection{Microscopic Origin of Diamagnetism (Begun)}\label{subsec:lec18-diamagnetism}

\subsubsection*{Physical Motivation}
Where do $\chi$ and its sign come from? We sketch the classical (Langevin) picture of
\emph{diamagnetism}, which arises from the \emph{orbital} motion of atomic electrons; the
spin contribution drives paramagnetism and is set aside here. Consider a gas of atoms,
each carrying orbital angular momentum $\bvec{L}$.

\subsubsection*{Larmor Precession}
Switching on a field $\bB$ (uniform on atomic scales) exerts a torque on the atomic moment
$\bmu=\dfrac{e}{2m_ec}\bvec{L}$, so
\begin{equation}\label{eq:lec18-torque}
  \frac{d\bvec{L}}{dt}=\bmu\times\bB
  =\frac{e}{2m_ec}\,\bvec{L}\times\bB
  =\boldsymbol{\omega}\times\bvec{L},
\end{equation}
which is the equation of \emph{precession} of $\bvec{L}$ about $\bB$ at the Larmor
frequency
\begin{equation}\label{eq:lec18-larmor}
  \boxed{\boldsymbol{\omega}=-\frac{e}{2m_ec}\,\bB.}
\end{equation}
\textit{Significance:} Every atomic moment precesses rigidly about $\bB$ at a single,
field-proportional rate $\omega$; since the electron charge is negative, $\boldsymbol{\omega}$
points along $\bB$.

\subsubsection*{Induced Moment}
The precession adds, on top of the orbital motion, a common circulation: each electron
$n$ acquires an extra velocity
\begin{equation}\label{eq:lec18-deltav}
  \Delta\bvec{v}_n=\boldsymbol{\omega}\times\br_n,
\end{equation}
hence an extra current density and an extra magnetic moment. With
$\Delta\Jvec=\sum_n e\,\Delta\bvec{v}_n\,\delta^{(3)}(\br-\br_n)$ and
$\Delta\bmu=\tfrac{1}{2c}\int \br\times\Delta\Jvec\,d^3r$, the $\delta$-functions collapse
the integral to
\begin{equation}\label{eq:lec18-deltamu}
  \Delta\bmu=\frac{e}{2c}\sum_{n}\br_n\times\bigl(\boldsymbol{\omega}\times\br_n\bigr).
\end{equation}
\textit{Significance:} The induced moment is quadratic in the geometry and (through
$\boldsymbol{\omega}\propto\bB$) linear in $\bB$, and it opposes $\bB$ -- the seed of a
\emph{negative} susceptibility. Completing the average over the precession and summing
over atoms per unit volume to extract $\chi$ is left for next lecture.

\nt{This derivation needs the \emph{orbital} angular momentum: spin also precesses, but
spin is not the circulation of charge, so it does not generate the velocity
\eqref{eq:lec18-deltav}. Spin instead underlies paramagnetism. The next lecture averages
$\Delta\bmu$ over the precession, sums over atoms, reads off $\chi$, and estimates its
(tiny) magnitude.}

%------------------------------------------------------------
\subsection{Summary}\label{subsec:lec18-summary}

\begin{itemize}
  \item The curl form of Gauss's theorem,
        $\int_V\nabla\times\bvec{A}\,d^3r=\oint_S\uvec{n}\times\bvec{A}\,d\sigma$, follows
        from the divergence theorem by dotting with a constant vector.
  \item $\psi_M$ is the electrostatic potential of effective magnetic charges
        $\rho_M=-\divg\bvec{M}$ (volume) and $\sigma_M=\bvec{M}\cdot\uvec{n}$ (surface);
        a uniform magnet has charges only on the caps perpendicular to $\bvec{M}$.
  \item Matching conditions: $B^{\perp}$ always continuous; $H^{\parallel}$ continuous
        unless a free surface current $\bvec{K}_{\rm f}$ flows.
  \item Linear media: $\bvec{M}=\chi\Hvec$, $\bB=\mu\Hvec$, $\mu=1+4\pi\chi$;
        $\chi>0$ paramagnet, $\chi<0$ diamagnet, $\chi=0$ vacuum.
  \item The magnetostatic BVP for $\mu>0$ has a unique solution
        ($\int\mu|\Hvec'|^2=0\Rightarrow\Hvec'=0$); consequently linear media keep no
        magnetization without external currents, unlike ferromagnets.
  \item Force on linear matter $\bvec{f}=\nabla \!\bigl(\tfrac{\chi}{2\mu}B^2\bigr)$:
        paramagnets are pulled into strong field, diamagnets expelled to weak field.
  \item Diamagnetism originates in Larmor precession
        $\boldsymbol{\omega}=-\tfrac{e}{2m_ec}\bB$ of orbital electrons, giving an induced
        moment $\Delta\bmu=\tfrac{e}{2c}\sum_n\br_n\times(\boldsymbol{\omega}\times\br_n)$
        opposing $\bB$ -- to be completed next lecture.
\end{itemize}
